%%
%% O5_B3_airy_gap_coercivity.tex
%%
%% Route B für O5: inf sigma(D^Airy) >= c^{1/3}/log(c) - ||K^Airy|| > 0
%%
%% DAG-Position:
%%   requires: O5-B-1  (||K^Airy|| <= kappa_BB, BB-Import Tracy-Widom)
%%             O5-B-2.1 [O]  (Airy-Fenster-Primzahldichte — SHORT INTERVAL OPEN)
%%             O5-B-2.2 [M]  (Frame-Schranke, importiert O4-B-2, unter 2C*zeta(3/2)<1)
%%             O5-B-2.3 [M]  (lambda_min(A^arith) >= c_frame)
%%             O4-B-2  [✓]   (Off-diagonal Airy-Abfall)
%%   closes:   O5 / P12  (bedingt auf O5-B-2.1 [O])
%%   feeds:    P13 (Main Theorem)
%%
%% STATUSMARKER-LEGENDE:
%%   [M]  = Mathematisch vollständig, bereit für Referee
%%   [O]  = Offen, noch nicht beweisbar ohne Zusatzannahme
%%   [H]  = Bedingt unter RH oder Cramer-Vermutung
%%   [BB] = Black-Box-Import aus Literatur, Referenz angegeben
%%   [?]  = Numerischer Wert noch nicht exakt rückgebunden
%%
%% Stand: 11. Juli 2026
%%

\documentclass[12pt,a4paper]{amsart}
\usepackage{amsmath,amssymb,amsthm,mathtools,enumitem,booktabs,hyperref,xcolor}

%% Prüfmarker-Makros
\newcommand{\statusM}{{\color{green!60!black}[\textbf{M}]}}
\newcommand{\statusO}{{\color{red}[\textbf{O}]}}
\newcommand{\statusH}{{\color{orange}[\textbf{H}]}}
\newcommand{\statusBB}{{\color{blue}[\textbf{BB}]}}
\newcommand{\statusQ}{{\color{red}[\textbf{?}]}}
\newcommand{\roteFlagge}[1]{{\color{red}\textbf{ROTE FLAGGE: }#1}}

\newtheorem{theorem}{Theorem}[section]
\newtheorem{lemma}[theorem]{Lemma}
\newtheorem{proposition}[theorem]{Proposition}
\newtheorem{corollary}[theorem]{Corollary}
\theoremstyle{definition}
\newtheorem{definition}[theorem]{Definition}
\newtheorem{assumption}[theorem]{Assumption}
\newtheorem{remark}[theorem]{Remark}

%% Notation
\newcommand{\KAiry}{K^{\mathrm{Airy}}}
\newcommand{\Aarith}{A^{\mathrm{arith}}}
\newcommand{\DAiry}{D^{\mathrm{Airy}}}
\newcommand{\PiAiry}{\Pi_{\mathrm{Airy}}}
\newcommand{\FTW}{F_2}
\newcommand{\kappaBB}{\kappa_{\mathrm{BB}}}
\newcommand{\cframe}{c_{\mathrm{frame}}}
\newcommand{\abs}[1]{\lvert #1 \rvert}
\newcommand{\norm}[1]{\lVert #1 \rVert}
\newcommand{\ip}[2]{\langle #1, #2 \rangle}

\begin{document}

\title{Lemma O5-B-3: Airy Gap Coercivity\\
$\inf\sigma(\DAiry) \ge \cframe(c) - \kappaBB > 0$\\
{\normalsize \statusO\ Bedingt auf O5-B-2.1 (Short-Interval-Primzahldichte)}}
\date{11.\ Juli 2026}
\maketitle

\begin{abstract}
Wir zeigen, dass $\inf\sigma(\DAiry) > 0$ gilt, sofern drei Voraussetzungen
erfüllt sind: (i)~eine numerisch rückzubindende Operatornorm-Schranke
$\norm{\KAiry} \le \kappaBB < 1$ (O5-B-1~\statusBB\statusQ),
(ii)~eine Primzahldichte auf Airy-Skala $c^{1/3}$, die \emph{nicht}
aus dem Standard-PNT folgt (O5-B-2.1~\statusO), und
(iii)~die explizite Positivitätsbedingung $2C\zeta(3/2) < 1$
(O5-B-2.2~\statusM\ unter Assumption~\ref{ass:schur}).
Der Hauptsatz O5-B-3 ist \statusM\ \emph{relativ zu} O5-B-2.1/B-2.2;
der Gesamtstatus von O5 bleibt \statusO.
\end{abstract}

\tableofcontents

%%=================================================================
\section{Setup und Hauptaussage}
\label{sec:setup}
%%=================================================================

\subsection{Der Operator}

\begin{definition}
Der \emph{arithmetische Mismatch-Operator im Airy-Fenster} ist
\[
  \DAiry := \PiAiry\bigl(\Aarith - \KAiry\bigr)\PiAiry,
\]
wobei
\begin{itemize}
  \item $\Aarith f = (1/N)\sum_j f(p_j^{\mathrm{Airy}})\,\delta_{p_j^{\mathrm{Airy}}}$
    mit $p_j^{\mathrm{Airy}} = c^{2/3}(p_j^* - t_N)$;
  \item $\KAiry$ der Airy-Kern-Operator auf $L^2([0,R])$
    (Tracy--Widom GUE, Standard-GUE-Normierung).
\end{itemize}
\end{definition}

\subsection{Hauptsatz}

\begin{theorem}[O5-B-3 \statusM\ relativ zu O5-B-2.1/B-2.2]
\label{thm:O5B3}
Unter den Voraussetzungen von O5-B-1 (Assumption~\ref{ass:kappaBB}),
O5-B-2.1 (Assumption~\ref{ass:density}) und O5-B-2.2
(Assumption~\ref{ass:schur}) gilt:
\[
  \inf\sigma(\DAiry)
  \;\ge\;
  \cframe(c)\bigl(1 - 2C\,\zeta\!\left(\tfrac{3}{2}\right)\bigr)
  - \kappaBB
  \;\xrightarrow{c\to\infty}\; +\infty.
\]
Insbesondere existiert ein explizites $c_0 > 0$ derart, dass
$\inf\sigma(\DAiry) > 0$ für alle $c \ge c_0$.
\end{theorem}

\begin{remark}[Gesamtstatus]
Der Satz ist logisch vollständig \statusM, aber sein Gesamtstatus
ist \statusO, da O5-B-2.1 offen ist; siehe Abschnitt~\ref{sec:O5B21}.
\end{remark}

%%=================================================================
\section{Lemma O5-B-1: Operatornorm von $\KAiry$}
\label{sec:O5B1}
%%=================================================================

%% ----------------------------------------------------------------
%% ROTE FLAGGE 1: numerischer Wert noch nicht rückgebunden
%% ----------------------------------------------------------------
\begin{remark}[\roteFlagge{Numerischer Wert von $\kappaBB$} \statusQ]
\label{rem:kappaBB_rot}
Die im Angriffs-Plan verwendete Zahl $1 - \FTW(0) \approx 0.0403$
entspricht \emph{nicht} dem Standardwert der GUE Tracy--Widom-Verteilung.
Aus Bornemann (2010) und Tracy--Widom (1994) ist:
\[
  F_2(0) \approx 0.9231, \qquad 1 - F_2(0) \approx 0.0769.
\]
Der Wert $0.9597$ (bzw.\ $1-0.9597 = 0.0403$) muss noch an eine
präzise Normierungskonvention für $\KAiry$ auf dem \emph{endlichen}
Intervall $[0,R]$ rückgebunden werden.
Bis dahin schreiben wir:
\[
  \kappaBB := 1 - \FTW(0;\text{Normierung}) < 1,
\]
ohne einen Dezimalwert einzusetzen.
\end{remark}

\begin{assumption}[O5-B-1 \statusBB\statusQ]
\label{ass:kappaBB}
Es existiert eine Konstante $\kappaBB \in (0,1)$, die aus dem
GUE Tracy--Widom Fredholm-Determinanten
$\det(I - \KAiry|_{[0,R]}) = \FTW(s_0)$
bei geeignetem $s_0$ (abhängig von der Intervall-Normierung) gewonnen wird,
derart dass
\[
  \norm{\KAiry}_{L^2([0,R])} \le \kappaBB < 1.
\]
\end{assumption}

\begin{lemma}[O5-B-1 \statusBB]
\label{lem:O5B1}
Unter Assumption~\ref{ass:kappaBB} gilt: für einen positiven Spurklasse-Operator
$K$ mit Eigenwerten $0 \le \lambda_j < 1$ und
$\det(I-K) = \prod_j(1-\lambda_j)$ folgt
\[
  \lambda_1 = \norm{K} \le 1 - \det(I-K),
\]
da $\prod_j(1-\lambda_j) \le 1 - \lambda_1$.
Angewendet auf $\KAiry$: $\norm{\KAiry} \le \kappaBB$.
\end{lemma}

\begin{proof}
Für $\lambda_j \in [0,1)$ paarweise: $\prod_{j\ge 1}(1-\lambda_j) \le 1-\lambda_1$,
denn $\prod_{j\ge 2}(1-\lambda_j) \le 1$. Daher $1-\lambda_1 \ge \det(I-K)$,
also $\lambda_1 \le 1 - \det(I-K) = \kappaBB$.
\end{proof}

\begin{remark}[Exaktwert \statusQ]
\label{rem:exact_value}
Sobald die Airy-Intervall-Normierung ($R$, Skalierung) exakt festgelegt ist,
muss $\kappaBB$ numerisch gegen Bornemann (2010, Math.\ Comp.) verifiziert werden.
Bis dahin bleibt O5-B-1 im Status \statusBB\statusQ.
\end{remark}

%%=================================================================
\section{Lemma O5-B-2: Arithmetische Sampling-Koerzivität}
\label{sec:O5B2}
%%=================================================================

\subsection{O5-B-2.1: Airy-Fenster-Primzahldichte \statusO}
\label{sec:O5B21}

%% ----------------------------------------------------------------
%% ROTE FLAGGE 2: Short-interval-Problem, nicht aus PNT
%% ----------------------------------------------------------------
\begin{remark}[\roteFlagge{Short-Interval-Problem} \statusO]
\label{rem:short_interval}
Die folgende Aussage über die Primzahldichte im Airy-Fenster erfordert
Primzahlen in einem Intervall der Länge $\sim c^{1/3}$ um $p_N \sim c\log c$.
Dies ist \emph{kein} Standard-PNT-Ergebnis:
\begin{itemize}
  \item Baker--Harman--Pintz (2001) und Verbesserungen (2023, Exponent $\approx 0.525$)
    liefern Primzahlen in $[x - x^{0.525}, x]$ unbedingt.
  \item Aber $x^{1/3} \ll x^{0.525}$: das Airy-Fenster ist
    \emph{erheblich kürzer} als der beste bekannte unbedingte Exponent.
  \item Unter RH gilt Primzahlen in $[x - x^{1/2+\varepsilon}, x]$,
    immer noch größer als $x^{1/3}$.
  \item Unter der Cramér-Vermutung (Intervalle $\sim \log^2 x$) wäre
    es beweisbar, aber das ist eine sehr starke Annahme.
\end{itemize}
\textbf{Konsequenz:} O5-B-2.1 wird als \statusO\ markiert.
Zwei saubere Alternativen:
\begin{enumerate}
  \item \statusH: Ergebnis unter RH formulieren
    ($[x-x^{1/2+\varepsilon}, x]$ enthält $\sim x^{1/2+\varepsilon}/\log x$ Primzahlen).
  \item Reformulierung auf größere Skala ($\sim c/\log c$, PNT-kompatibel),
    dann Airy-Fenster via Projektion/Mittelung, ohne direktes Short-interval-Zählen.
\end{enumerate}
\end{remark}

\begin{assumption}[O5-B-2.1 \statusO]
\label{ass:density}
Die Airy-reskalierten Primzahlen $\{p_j^{\mathrm{Airy}}\}$ im Fenster $[0,R_c]$
mit $R_c = c^{2/3}\delta$ erfüllen:
\[
  \#\{j : p_j^{\mathrm{Airy}} \in [0,R_c]\}
  \;\asymp\; \frac{c^{1/3}}{\log c}.
\]
\textbf{Status \statusO:}
Dies ist äquivalent zur Existenz $\sim c^{1/3}/\log c$ vieler Primzahlen
in einem Intervall der Länge $\sim c^{1/3}$ um $c\log c$,
was den besten unbedingten Short-interval-Resultaten überlegen ist.
\end{assumption}

\begin{lemma}[O5-B-2.1 \statusO, bedingt]
\label{lem:O5B21}
Unter Assumption~\ref{ass:density} konvergiert die empirische Dichte
der $\{p_j^{\mathrm{Airy}}\}$ auf kompakten Teilmengen von $[0,\infty)$
zu einem absolutstetigen Maß mit Dichte
\[
  \rho_{\mathrm{Airy}}(t) \;=\; \frac{A_\delta}{\log c}\cdot c^{1/3}(1+o(1)),
  \qquad A_\delta > 0 \text{ explizit aus PNT}.
\]
\end{lemma}

\begin{proof}
Unter Assumption~\ref{ass:density}: Es gibt $\sim c^{1/3}/\log c$ Punkte
in einem Intervall der (Airy-reskalierten) Länge $R_c \sim c^{2/3}$,
also Dichte $\sim (c^{1/3}/\log c) / c^{2/3} = c^{-1/3}/\log c$ — dies
ist die normierte Dichte pro Einheitslänge in Airy-Koordinaten.
Die absolute Stetigkeit und Gleichmäßigkeit auf Kompakta
folgt aus der Gleichverteilung der Primzahlen (PNT mit Fehlerterm,
Dusart 2010), sobald Assumption~\ref{ass:density} gilt.
\end{proof}

\subsection{O5-B-2.2: Frame-Schranke in Airy-Koordinaten \statusM}
\label{sec:O5B22}

%% ----------------------------------------------------------------
%% Positivitätsbedingung explizit
%% ----------------------------------------------------------------
\begin{assumption}[Schur-Positivität \statusM]
\label{ass:schur}
Die Konstante $C > 0$ aus dem Off-Diagonal-Abfall (O4-B-2) erfüllt:
\[
  2C\,\zeta\!\left(\tfrac{3}{2}\right) < 1.
\]
Nur unter dieser Bedingung ist die untere Frame-Schranke in Lemma~\ref{lem:O5B22}
positiv.
\end{assumption}

\begin{remark}
$\zeta(3/2) \approx 2.612$, also ist Assumption~\ref{ass:schur}
äquivalent zu $C < 1/(2 \cdot 2.612) \approx 0.191$.
Ob dies für die tatsächliche Konstante $C$ aus O4-B-2 gilt,
ist numerisch zu überprüfen (vgl.\ Offene Frage~\ref{oq:C_value}).
\end{remark}

\begin{lemma}[O5-B-2.2 \statusM\ unter Assumptions~\ref{ass:density},\ref{ass:schur}]
\label{lem:O5B22}
Sei $\{\phi_k\}_{k\ge 0}$ eine $L^2$-ONB aus Airy-Eigenfunktionen.
Für $f = \sum_k a_k\phi_k \in \PiAiry L^2([0,R])$:
\[
  \frac{1}{N}\sum_{j=1}^N \abs{f(p_j^{\mathrm{Airy}})}^2
  \;\ge\;
  \cframe(c)\bigl(1 - 2C\,\zeta\!\left(\tfrac{3}{2}\right)\bigr)\norm{f}^2,
\]
wobei $\cframe(c) = A_\delta\,c^{1/3}/\log c > 0$
und die rechte Seite positiv ist unter Assumption~\ref{ass:schur}.
\end{lemma}

\begin{proof}
Aufspaltung in Diagonal- und Off-Diagonal-Terme:
\[
  \frac{1}{N}\sum_j \abs{f}^2
  = \sum_k \abs{a_k}^2 \underbrace{\frac{1}{N}\sum_j \abs{\phi_k(p_j^{\mathrm{Airy}})}^2}_{\ge\,\cframe(c)\text{ (O5-B-2.1)}}
  + \sum_{k\ne l} a_k\overline{a_l}\,H_{kl}.
\]
Die Off-Diagonal-Einträge erfüllen $\abs{H_{kl}} \le C/\abs{k-l}^{3/2}$
(Airy-Oszillationsabfall, Struktur identisch zu O4-B-2).
Schur-Test:
\[
  \abs{\sum_{k\ne l} a_k\overline{a_l}\,H_{kl}}
  \le \norm{f}^2 \sup_k \sum_{l\ne k}\frac{C}{\abs{k-l}^{3/2}}
  = \norm{f}^2 \cdot 2C\,\zeta\!\left(\tfrac{3}{2}\right).
\]
Kombination und Anwendung von Assumption~\ref{ass:schur}
ergibt die Behauptung.
\end{proof}

\subsection{O5-B-2.3: Spektrale Koerzivität von $\Aarith$ \statusM}
\label{sec:O5B23}

\begin{lemma}[O5-B-2.3 \statusM\ unter Assumptions~\ref{ass:density},\ref{ass:schur}]
\label{lem:O5B23}
\[
  \lambda_{\min}\bigl(\Aarith\big|_{\mathrm{edge}}\bigr)
  \;\ge\;
  \cframe(c)\bigl(1 - 2C\,\zeta\!\left(\tfrac{3}{2}\right)\bigr)
  \;=:\; \mu_{\min}(c)
  \;\sim\; \frac{c^{1/3}}{\log c}.
\]
\end{lemma}

\begin{proof}
Rayleigh-Quotient-Charakterisierung und Lemma~\ref{lem:O5B22}.
\end{proof}

%%=================================================================
\section{Beweis von Theorem~\ref{thm:O5B3}}
\label{sec:proof}
%%=================================================================

\begin{proof}[Beweis von Theorem~\ref{thm:O5B3}]
Min-Max-Prinzip für selbstadjungierte Operatoren:
\begin{align}
  \inf\sigma(\DAiry)
  &= \inf\sigma\bigl(\PiAiry(\Aarith - \KAiry)\PiAiry\bigr) \notag\\
  &\ge \lambda_{\min}(\Aarith|_{\mathrm{edge}})
       - \norm{\KAiry|_{\mathrm{edge}}} \notag\\
  &\ge \mu_{\min}(c) - \kappaBB
  \;\ge\; \frac{A_\delta(1-2C\zeta(3/2))}{\log c}\cdot c^{1/3} - \kappaBB.
  \label{eq:hauptsatz}
\end{align}
Da $c^{1/3}/\log c \to \infty$ und $\kappaBB < 1$ konstant:
existiert $c_0$ explizit aus PNT derart, dass die rechte Seite
für $c \ge c_0$ positiv ist.
\end{proof}

%%=================================================================
\section{DAG-Abschluss}
\label{sec:dag}
%%=================================================================

\begin{center}
\renewcommand{\arraystretch}{1.4}
\begin{tabular}{lllll}
\toprule
\textbf{Node} & \textbf{Aussage (kurz)} & \textbf{Status} & \textbf{Typ} \\
\midrule
O5-B-1    & $\norm{\KAiry} \le \kappaBB$                     & \statusBB\statusQ & Import \\
O5-B-2.1  & Airy-Fenster-Dichte $\asymp c^{1/3}/\log c$     & \statusO         & Offen  \\
O5-B-2.2  & Frame-Schranke unter $2C\zeta(3/2)<1$            & \statusM         & Beweis \\
O5-B-2.3  & $\lambda_{\min}(\Aarith) \ge \mu_{\min}(c)$     & \statusM         & Beweis \\
O5-B-3    & $\inf\sigma(\DAiry) > 0$                         & \statusM\ rel.\ B-2.1/2.2 & Hauptsatz \\
\midrule
O5 / P12  & $\inf\sigma(\DAiry) \ge c_2 > 0$                 & \statusO         & Ziel   \\
P13       & Main Theorem (O4+O5)                              & bedingt \statusO  & Ziel   \\
\bottomrule
\end{tabular}
\end{center}

\medskip
\noindent\textbf{Kritischer Pfad:}
\[
  \underbrace{\text{O4}}_{\checkmark}
  \Rightarrow P9 \Rightarrow \text{R1}
  \qquad
  \underbrace{\text{O5-B-2.1}}_{\statusO}
  \to \text{O5-B-2.2} \to \text{O5-B-3}
  \Rightarrow P12
  \;\Rightarrow\; \boxed{P13}.
\]
Das einzige verbleibende echte Hindernis im Minimalpfad ist O5-B-2.1.

%%=================================================================
\section{Offene Fragen}
\label{sec:open}
%%=================================================================

\begin{enumerate}[label={[O-O5B3-\arabic*]},ref={O-O5B3-\arabic*}]

\item\label{oq:kappaBB_exact}
\textbf{Exakter Wert von $\kappaBB$.}
Welche Normierungskonvention für $\KAiry$ auf $[0,R]$ liefert
$\kappaBB \approx 0.0403$? Oder ist dies ein Fehler im Angriffs-Plan?
Bornemann (2010, Math.\ Comp.) numerisch prüfen.
\statusQ

\item\label{oq:O5B21_resolution}
\textbf{Auflösung von O5-B-2.1.}
Drei Optionen:
(a)~Bedingung \statusH\ unter RH ($[x-x^{1/2+\varepsilon},x]$-Primzahlen);
(b)~Reformulierung auf größere Skala ($c/\log c$), Airy-Fenster via Projektion;
(c)~Als \statusO\ belassen und P12 als bedingt markieren.
\statusO

\item\label{oq:C_value}
\textbf{Wert der Off-Diagonal-Konstante $C$.}
Ist $C < 0.191$ (Assumption~\ref{ass:schur} erfüllt)?
Numerisch aus der Airy-Eigenfunktionsbasis abschätzen.
\statusQ

\end{enumerate}

\begin{thebibliography}{99}
\bibitem{TracyWidom1994}
C.~A.~Tracy und H.~Widom,
\textit{Level-spacing distributions and the Airy kernel},
Comm.\ Math.\ Phys.\ \textbf{159} (1994), 151--174.

\bibitem{Bornemann2010}
F.~Bornemann,
\textit{On the numerical evaluation of distributions in random matrix theory:
a review},
Markov Process.\ Related Fields \textbf{16} (2010), 803--866.

\bibitem{BakerHarmanPintz2001}
R.~C.~Baker, G.~Harman und J.~Pintz,
\textit{The difference between consecutive primes, II},
Proc.\ London Math.\ Soc.\ \textbf{83} (2001), 532--562.

\bibitem{Dusart2010}
P.~Dusart,
\textit{Estimates of some functions over primes without $R$.\,H.},
Preprint (2010), arXiv:1002.0442.

\bibitem{O4B3}
S.~Schmalnauer,
\texttt{O4\_B3\_jaffard\_invertibility.tex},
\textsf{prolate-gram-coercivity}, 2026.
\end{thebibliography}

\end{document}
